% ৫. মধ্যস্তরের ভাষা
\section{মধ্যস্তরের ভাষা}

\begin{learningbox}
এই টপিক শেষে শিক্ষার্থী mid-level language কেন বলা হয় এবং C ভাষার ভূমিকা ব্যাখ্যা করতে পারবে।
\end{learningbox}

\begin{definitionbox}{মধ্যস্তরের ভাষা}
যে ভাষায় high-level language-এর মতো সহজ code লেখা যায়, আবার low-level operation বা hardware control-ও করা যায়, তাকে mid-level language বলা হয়।
\end{definitionbox}

\begin{examplebox}{C কেন mid-level}
C ভাষায় function, loop, data type দিয়ে structured program লেখা যায়; আবার pointer, bitwise operator ও memory-level কাজও করা যায়। তাই C-কে অনেক সময় mid-level language বলা হয়।
\end{examplebox}

\begin{boardbox}{বোর্ড ফোকাস}
C ভাষাকে mid-level বলার কারণ লিখতে high-level readability এবং low-level hardware control—দুই দিকই উল্লেখ করো।
\end{boardbox}
